% F2.2 -- Hand proof that every constraint of the reduced O(K^2) LP
% (code/reduced_lp.py, RESEARCH_STATE R6) is a valid inequality for every
% instance (f, ftilde) and the predictive-greedy trajectory on it.
%
% STATUS: [HAND-PROOF-UNREVIEWED].  No oracle certifies the derivations below.
% What IS oracle-checked (see RESEARCH_STATE R6, R10 and results/N1, N3):
%   - the reduced LP value agrees with the full-lattice LP at 19 test points
%     and for K <= 5  [VERIFIED-LP];
%   - the closed form min_j V_j of the reduced LP has a symbolic dual
%     certificate for general K  [VERIFIED-SYMBOLIC].
% The present file supplies the missing direction: that the LP is a RELAXATION,
% which is what makes its optimum a lower bound on the true worst case for
% every K.  In particular the case b_t in O, which the earlier note flagged as
% the gap in the argument, is treated explicitly in every one of the four
% families (Lemmas 3 to 6, case (c) of each).
%
% This file is an \input fragment.  Host document must load amsmath and amsthm
% and must define the environments `lemma`, `definition`, `remark` and `proof`.
% Compile test: paper/test_F2.tex under the ICLR 2027 template, log in
% results/F2_compile.log.

\section{Validity of the reduced LP inequalities}
\label{sec:reduced-lp-validity}

This section verifies that each of the four constraint families of the reduced
$O(K^2)$ linear program is a valid inequality: for every instance and every run
of predictive greedy on it, the quantities read off the run satisfy the
constraint. Consequently the LP is a relaxation of the set of realisable runs,
its optimum is a lower bound on the worst-case ratio for every $K$, and the
closed form of that optimum inherits the same status.

\subsection{Setting and notation}

Let $N$ be a finite ground set and let $f \colon 2^{N} \to \mathbb{R}$ be
monotone and submodular with $f(\emptyset) = 0$. Write
$d_{e}(S) = f(S \cup \{e\}) - f(S)$ for the marginal gain of $e \notin S$.
The predictor $\tilde f \colon 2^{N} \to \mathbb{R}$ is an arbitrary set
function with $\tilde f(\emptyset) = 0$; it is not assumed monotone or
submodular. Its gains are written $\tilde d_{e}(S)$.

\begin{definition}[error band]
\label{def:band}
For $\etau, \etao \ge 1$ we say that $\tilde f$ is $(\etau, \etao)$-accurate if
for every $S \subseteq N$ and every $e \notin S$
\begin{equation}
\label{eq:band}
\frac{1}{\etau}\, d_{e}(S) \;\le\; \tilde d_{e}(S) \;\le\; \etao\, d_{e}(S) .
\end{equation}
We call the right inequality of \eqref{eq:band} the \emph{upper band} and the
left inequality the \emph{lower band}, following the constraint names
\texttt{band\_up} and \texttt{band\_lo} in the LP builders. The global error is
$\eta = \etau \etao \ge 1$. Gains are assumed strictly positive on the states
that occur below, which is the regime in which \eqref{eq:band} has content.
\end{definition}

\emph{Predictive greedy} starts from $S^{0} = \emptyset$ and for
$t = 0, \dots, K-1$ picks
\begin{equation}
\label{eq:greedy}
b_{t} \in \arg\max_{e \notin S^{t}} \tilde d_{e}(S^{t}),
\qquad S^{t+1} = S^{t} \cup \{b_{t}\},
\end{equation}
and returns $S^{K}$. Ties in \eqref{eq:greedy} are broken adversarially. We
refer to \eqref{eq:greedy} as the \emph{greedy choice}.

Fix an optimal set $O$ with $|O| \le K$ and $f(O) = \max_{|T| \le K} f(T)$.
Enumerate $O = \{o_{1}, \dots, o_{K}\}$ and normalise $f(O) = 1$. If $|O| < K$,
set $g_{t,i} = 0$ for the $K - |O|$ unused indices $i$; every lemma below then
holds for those columns for the same reason as in its case (a), so no padding of
the ground set is needed.

\begin{definition}[the trajectory vector, including the case $b_{t} \in O$]
\label{def:vars}
For $t = 0, \dots, K-1$ put $d_{t} = d_{b_{t}}(S^{t})$, and for
$t = 0, \dots, K$ and $i = 1, \dots, K$ put
\begin{equation}
\label{eq:gdef}
g_{t,i} \;=\;
\begin{cases}
d_{o_{i}}(S^{t}), & o_{i} \notin S^{t},\\[2pt]
0, & o_{i} \in S^{t}.
\end{cases}
\end{equation}
\end{definition}

The second line of \eqref{eq:gdef} is the \emph{zeroing convention}. It is the
device that handles the case in which greedy selects an element of $O$: as soon
as $b_{t} = o_{i}$, the variable $g_{t+1,i}$ and all later ones are set to zero
rather than left undefined. Every lemma below is stated and proved with this
convention in force, and in each proof the case $b_{t} = o_{i}$ appears as a
separate case (c). Note that the convention is consistent with reading
$g_{t,i}$ as a marginal gain, because an element already in $S^{t}$ contributes
no further gain to it.

Also set $r_{t} = 1 - \sum_{s<t} d_{s}$, the residual after $t$ steps. Because
$\sum_{s<t} d_{s}$ telescopes to $f(S^{t}) - f(\emptyset) = f(S^{t})$, we have
\begin{equation}
\label{eq:residual}
r_{t} \;=\; f(O) - f(S^{t}).
\end{equation}

The reduced LP of \texttt{code/reduced\_lp.py} is, in these variables,
\begin{equation}
\label{eq:redlp}
\begin{aligned}
\min \quad & \textstyle\sum_{t=0}^{K-1} d_{t}\\
\text{s.t.} \quad
& \textstyle\sum_{i=1}^{K} g_{t,i} \ \ge\ r_{t}
&& (\mathrm{sum}(t)), \quad t = 0,\dots,K-1,\\
& d_{t} \ \ge\ g_{t,i}/\eta
&& (\mathrm{pred}(t,i)),\\
& g_{t+1,i} \ \le\ g_{t,i}
&& (\mathrm{mono}(t,i)),\\
& (1 - 1/\eta)\, g_{t+1,i} \ \ge\ g_{t,i} - d_{t}
&& (\mathrm{cons}(t,i)),\\
& d_{t} \ \ge\ 0, \quad g_{t,i} \ \ge\ 0 .
\end{aligned}
\end{equation}

\subsection{The coherence lemma}

The fourth family rests on the coherence lemma (RESEARCH\_STATE R3, called the
consistency lemma in earlier drafts). We restate it and give its proof, so that
the present section is self contained.

\begin{lemma}[coherence, R3]
\label{lem:coherence}
Let $S \subseteq N$ and let $e, e' \notin S$ with $e \ne e'$ satisfy
$\tilde d_{e}(S) \ge \tilde d_{e'}(S)$. Then
\begin{align}
d_{e}(S \cup \{e'\}) &\ \ge\ d_{e'}(S \cup \{e\}) / \eta,
\label{eq:coh-i}\\
(1 - 1/\eta)\, d_{e'}(S \cup \{e\}) &\ \ge\ d_{e'}(S) - d_{e}(S).
\label{eq:coh-ii}
\end{align}
\end{lemma}

\begin{proof}
Expanding $\tilde f(S \cup \{e,e'\}) - \tilde f(S)$ in the two possible orders
gives the identity
$\tilde d_{e}(S) + \tilde d_{e'}(S \cup \{e\})
 = \tilde d_{e'}(S) + \tilde d_{e}(S \cup \{e'\})$.
This step uses nothing beyond the definition of a set function.

Subtracting and using $\tilde d_{e}(S) \ge \tilde d_{e'}(S)$ yields
$\tilde d_{e'}(S \cup \{e\}) \le \tilde d_{e}(S \cup \{e'\})$.
This step uses the greedy choice, in the form of the hypothesis on $S$.

Now $d_{e'}(S \cup \{e\}) \le \etau\, \tilde d_{e'}(S \cup \{e\})$; this step
uses the lower band. Then
$\etau\, \tilde d_{e'}(S \cup \{e\}) \le \etau\, \tilde d_{e}(S \cup \{e'\})$
by the previous display, and
$\etau\, \tilde d_{e}(S \cup \{e'\}) \le \etau \etao\, d_{e}(S \cup \{e'\})
 = \eta\, d_{e}(S \cup \{e'\})$; this step uses the upper band. Chaining the
three gives \eqref{eq:coh-i}.

For \eqref{eq:coh-ii}, the same two-order expansion applied to $f$ gives
$d_{e}(S) + d_{e'}(S \cup \{e\}) = d_{e'}(S) + d_{e}(S \cup \{e'\})$, hence
$d_{e'}(S) - d_{e}(S) = d_{e'}(S \cup \{e\}) - d_{e}(S \cup \{e'\})$. This step
uses nothing beyond the definition of a set function. Substituting
\eqref{eq:coh-i} in the form
$- d_{e}(S \cup \{e'\}) \le - d_{e'}(S \cup \{e\})/\eta$ gives
$d_{e'}(S) - d_{e}(S) \le (1 - 1/\eta)\, d_{e'}(S \cup \{e\})$, which is
\eqref{eq:coh-ii}. This step uses \eqref{eq:coh-i} together with $\eta \ge 1$,
so that the factor $1 - 1/\eta$ is nonnegative.
\end{proof}

\subsection{The four families are valid}

Throughout, $t \in \{0, \dots, K-1\}$ and $i \in \{1, \dots, K\}$, and the three
cases are
(a) $o_{i} \in S^{t}$;
(b) $o_{i} \notin S^{t}$ and $o_{i} \ne b_{t}$;
(c) $o_{i} = b_{t}$, which is the case $b_{t} \in O$.

\begin{lemma}[sign constraints]
\label{lem:sign}
$d_{t} \ge 0$ and $g_{t,i} \ge 0$ for all $t, i$.
\end{lemma}

\begin{proof}
Each of these numbers is either a marginal gain of $f$ or, under the zeroing
convention, equal to $0$. This step uses monotonicity of $f$.
\end{proof}

\begin{lemma}[$\mathrm{sum}(t)$: coverage]
\label{lem:sum}
For every $t$, $\ \sum_{i=1}^{K} g_{t,i} \ \ge\ r_{t}$.
\end{lemma}

\begin{proof}
Let $I_{t} = \{ i : o_{i} \notin S^{t} \}$, so that by the zeroing convention
$\sum_{i=1}^{K} g_{t,i} = \sum_{i \in I_{t}} d_{o_{i}}(S^{t})$. This step uses
Definition~\ref{def:vars}, and it is exactly where the case $b_{s} \in O$ for
some $s < t$ is absorbed: the indices of the elements of $O$ already picked
drop out of the sum, and they drop out on both sides of the argument below.

Order $I_{t} = \{ i_{1}, \dots, i_{m} \}$ and telescope
\begin{equation*}
f(S^{t} \cup O) - f(S^{t})
= \sum_{j=1}^{m} \Big[ f\big(S^{t} \cup \{o_{i_{1}}, \dots, o_{i_{j}}\}\big)
 - f\big(S^{t} \cup \{o_{i_{1}}, \dots, o_{i_{j-1}}\}\big) \Big]
\le \sum_{j=1}^{m} d_{o_{i_{j}}}(S^{t}).
\end{equation*}
This step uses submodularity, applied once per term: the gain of $o_{i_{j}}$ at
the larger set $S^{t} \cup \{o_{i_{1}}, \dots, o_{i_{j-1}}\}$ is at most its
gain at $S^{t}$.

Next, $f(S^{t} \cup O) \ge f(O)$. This step uses monotonicity.

Combining the two displays with \eqref{eq:residual},
\begin{equation*}
\sum_{i=1}^{K} g_{t,i} \ \ge\ f(S^{t} \cup O) - f(S^{t})
\ \ge\ f(O) - f(S^{t}) \ =\ r_{t}.
\qedhere
\end{equation*}
\end{proof}

\begin{lemma}[$\mathrm{pred}(t,i)$: greedy choice inside the band]
\label{lem:pred}
For every $t, i$, $\ d_{t} \ \ge\ g_{t,i} / \eta$.
\end{lemma}

\begin{proof}
Case (a): $g_{t,i} = 0$ and $d_{t} \ge 0$, so the inequality holds. This step
uses monotonicity (through Lemma~\ref{lem:sign}).

Case (b): chain
\begin{equation*}
d_{t} \;=\; d_{b_{t}}(S^{t})
\ \ge\ \frac{\tilde d_{b_{t}}(S^{t})}{\etao}
\ \ge\ \frac{\tilde d_{o_{i}}(S^{t})}{\etao}
\ \ge\ \frac{d_{o_{i}}(S^{t})}{\etau \etao}
\;=\; \frac{g_{t,i}}{\eta}.
\end{equation*}
The first inequality uses the upper band at $(S^{t}, b_{t})$, rearranged as
$d \ge \tilde d / \etao$. The second inequality uses the greedy choice, since
$o_{i} \notin S^{t}$ makes $o_{i}$ an eligible candidate at step $t$. The third
inequality uses the lower band at $(S^{t}, o_{i})$, rearranged as
$\tilde d \ge d / \etau$.

Case (c): here $g_{t,i} = d_{o_{i}}(S^{t}) = d_{b_{t}}(S^{t}) = d_{t}$, and the
claim reads $d_{t} \ge d_{t}/\eta$, which holds because $d_{t} \ge 0$ and
$\eta \ge 1$. This step uses monotonicity and $\eta \ge 1$. Note that no band
inequality is needed in this case, so the constraint does not degrade when
greedy happens to pick an optimal element.
\end{proof}

\begin{lemma}[$\mathrm{mono}(t,i)$: single-element submodularity]
\label{lem:mono}
For every $t, i$, $\ g_{t+1,i} \ \le\ g_{t,i}$.
\end{lemma}

\begin{proof}
Case (a): $o_{i} \in S^{t} \subseteq S^{t+1}$, so both sides are $0$. This step
uses the zeroing convention only.

Case (b): $o_{i} \notin S^{t+1}$ as well, so
$g_{t+1,i} = d_{o_{i}}(S^{t+1}) \le d_{o_{i}}(S^{t}) = g_{t,i}$ because
$S^{t} \subseteq S^{t+1}$. This step uses submodularity, in its
diminishing-returns form for a single element.

Case (c): $o_{i} = b_{t} \in S^{t+1}$, so $g_{t+1,i} = 0$ while
$g_{t,i} = d_{t} \ge 0$. This step uses monotonicity. This is the only place in
the four families where the case $b_{t} \in O$ replaces a submodularity argument
by a monotonicity argument.
\end{proof}

\begin{lemma}[$\mathrm{cons}(t,i)$: coherence]
\label{lem:cons}
For every $t, i$, $\ (1 - 1/\eta)\, g_{t+1,i} \ \ge\ g_{t,i} - d_{t}$.
\end{lemma}

\begin{proof}
Case (a): the left side is $0$ and the right side is $0 - d_{t} \le 0$. This
step uses monotonicity.

Case (b): apply Lemma~\ref{lem:coherence} with $S = S^{t}$, $e = b_{t}$ and
$e' = o_{i}$. The hypothesis $\tilde d_{b_{t}}(S^{t}) \ge \tilde d_{o_{i}}(S^{t})$
of that lemma holds because $o_{i}$ is eligible at step $t$; this step uses the
greedy choice. Inequality \eqref{eq:coh-ii} then reads
\begin{equation*}
(1 - 1/\eta)\, d_{o_{i}}(S^{t} \cup \{b_{t}\})
\ \ge\ d_{o_{i}}(S^{t}) - d_{b_{t}}(S^{t}),
\end{equation*}
which is the claim, since $S^{t+1} = S^{t} \cup \{b_{t}\}$ and $o_{i} \notin
S^{t+1}$ give $d_{o_{i}}(S^{t+1}) = g_{t+1,i}$. This step uses R3(ii).

Case (c): $g_{t+1,i} = 0$ by the zeroing convention and
$g_{t,i} - d_{t} = d_{b_{t}}(S^{t}) - d_{b_{t}}(S^{t}) = 0$, so the constraint
holds with equality. This step uses Definition~\ref{def:vars} only. In
particular Lemma~\ref{lem:coherence}, which requires $e \ne e'$, is never
invoked in this case, which is what closes the gap flagged in R6.
\end{proof}

\subsection{Consequence}

\begin{lemma}[relaxation]
\label{lem:relaxation}
For every instance $(f, \tilde f)$ satisfying Definition~\ref{def:band} with
$\eta = \etau \etao$, every optimal set $O$ and every run of predictive greedy,
the vector $(d_{t}, g_{t,i})$ of Definition~\ref{def:vars} is feasible for
\eqref{eq:redlp} and its objective value equals $f(S^{K}) / f(O)$. Hence the
optimum of \eqref{eq:redlp} is a lower bound on the worst-case ratio of
predictive greedy at budget $K$ and error $\eta$.
\end{lemma}

\begin{proof}
Feasibility is Lemmas~\ref{lem:sign} to \ref{lem:cons}. The objective is
$\sum_{t<K} d_{t} = f(S^{K})$ by telescoping, and $f(O) = 1$ by the
normalisation, so the objective equals the ratio. This step uses the definition
of $d_{t}$ and the normalisation. Since every realisable run gives a feasible
point, the minimum over the feasible set is at most the minimum over realisable
runs, that is, at most the worst-case ratio.
\end{proof}

\begin{remark}[what is and is not claimed]
\label{rem:status}
Lemma~\ref{lem:relaxation} gives one direction only. The reverse direction,
that every feasible point of \eqref{eq:redlp} is realised by some instance, is
not claimed here; it has been checked against the full-lattice LP for
$K \le 5$ and at nineteen $(K, \eta)$ test points. The status of the material in
this section is [HAND-PROOF-UNREVIEWED]: no oracle certifies these derivations.
\end{remark}

\begin{remark}[which rulers the four families tolerate]
\label{rem:rulers}
Families $\mathrm{sum}(t)$ and $\mathrm{mono}(t,i)$ use no error notion at all.
Family $\mathrm{pred}(t,i)$ uses the band only at the greedy state $S^{t}$, so
it stays valid when $\eta$ is replaced by the trajectory error $\etatr$, and,
directly from its definition as a ratio of true gains at greedy states, when it
is replaced by the selection error $\etasel$. Family $\mathrm{cons}(t,i)$
invokes Lemma~\ref{lem:coherence}, whose proof also uses the band at the state
$S^{t} \cup \{o_{i}\}$, which is not on the greedy trajectory; so
$\mathrm{cons}$ is not covered by either restricted ruler without a further
assumption. This is the reason the $L_{K}$ bound, whose proof uses only
$\mathrm{sum}$ and $\mathrm{pred}$, is stated for $\etasel$, while the exact
reduced-LP value is stated for the global $\eta$.
\end{remark}

\begin{remark}[which hypothesis each family consumes]
\label{rem:census}
Family $\mathrm{sum}(t)$ uses submodularity and monotonicity of $f$ and no
property of $\tilde f$. Family $\mathrm{pred}(t,i)$ uses the greedy choice and
both bands, and no property of $f$ beyond monotonicity. Family
$\mathrm{mono}(t,i)$ uses submodularity, and monotonicity in case (c). Family
$\mathrm{cons}(t,i)$ uses the greedy choice and both bands through
Lemma~\ref{lem:coherence}, and monotonicity in cases (a) and (c). No family
uses submodularity of $\tilde f$, which is consistent with the predictor being
an arbitrary set function.
\end{remark}
